% Independently written from Illusie I, 2.3.1, printed pp.38--42.
% Local labels only. Source snippet for bounded root-chapter composition.

\section{The topos of simplicial sheaves}
\label{section-illusie-I-simplicial-topos}

\begin{lemma}
\label{lemma-illusie-I-simplicial-sheaves-topos}
Let $\mathcal{C}$ be a site and put $\mathcal{T} = \Sh(\mathcal{C})$.
Then the category $\text{Simp}(\mathcal{T})$ of simplicial sheaves is a
topos.  Moreover, for $U \in \mathcal{T}$, $n \geq 0$, and
$X \in \text{Simp}(\mathcal{T})$ there is a canonical bijection
$$
\Mor_{\text{Simp}(\mathcal{T})}(U \times \Delta[n], X)
=
\Mor_{\mathcal{T}}(U, X_n).
$$
Consequently, if a family $\mathcal{G}$ detects equality of arrows in
$\mathcal{T}$, then the objects
$$
U \times \Delta[n], \qquad U \in \mathcal{G},\ n \geq 0,
$$
detect equality of arrows in $\text{Simp}(\mathcal{T})$.

For $X,Y \in \text{Simp}(\mathcal{T})$ the internal hom
$\SheafHom(X,Y)$ exists.  Its $n$th term is characterized by
$$
\Mor_{\mathcal{T}}(U,\SheafHom(X,Y)_n)
=
\Mor_{\text{Simp}(\mathcal{T})}
(U \times \Delta[n] \times X,Y),
$$
functorially in $U$, and there are functorial isomorphisms
$$
\SheafHom(Z,\SheafHom(X,Y))
\longrightarrow
\SheafHom(Z \times X,Y).
$$
\end{lemma}

\begin{proof}
Let $\Delta$ denote the simplex category.  On
$\mathcal{C} \times \Delta$ take the topology generated by the families
$$
(U_i,[n]) \longrightarrow (U,[n])
$$
for which $\{U_i \to U\}$ is a covering of $\mathcal{C}$.  A presheaf on
$\mathcal{C} \times \Delta$ is the same thing as a simplicial presheaf on
$\mathcal{C}$.  The sheaf condition for the displayed generating families
says exactly that every degree is a sheaf on $\mathcal{C}$.  We therefore
have an equivalence
$$
\Sh(\mathcal{C} \times \Delta)
\longrightarrow
\text{Simp}(\Sh(\mathcal{C})).
$$
Thus $\text{Simp}(\mathcal{T})$ is a topos by Sites, Definition
\ref{sites-definition-topos}.

The asserted bijection is Simplicial Methods, Lemma
\ref{simplicial-lemma-morphism-from-coproduct}.  It also shows the assertion
about a detecting family: equality in every degree is equality of simplicial
arrows.

Since $\text{Simp}(\mathcal{T})$ is a category of sheaves on a site, the
internal hom exists by Sites, Lemma
\ref{sites-lemma-internal-hom-sheaf}.  Applying the preceding bijection and
then the defining adjunction gives
\begin{align*}
\Mor_{\mathcal{T}}(U,\SheafHom(X,Y)_n)
&=
\Mor_{\text{Simp}(\mathcal{T})}
(U \times \Delta[n],\SheafHom(X,Y)) \\
&=
\Mor_{\text{Simp}(\mathcal{T})}
(U \times \Delta[n] \times X,Y).
\end{align*}
The last displayed isomorphism in the statement follows by applying the
same adjunction twice and using the Yoneda lemma.
\end{proof}

\begin{lemma}
\label{lemma-illusie-I-simplicial-topos-morphism}
Let
$$
f : \Sh(\mathcal{D}) \longrightarrow \Sh(\mathcal{C})
$$
be a morphism of topoi.  Applying $f^{-1}$ and $f_*$ degreewise defines a
morphism of topoi
$$
\text{Simp}(f) :
\text{Simp}(\Sh(\mathcal{D}))
\longrightarrow
\text{Simp}(\Sh(\mathcal{C})).
$$
For simplicial sheaves $X,Y$ on $\mathcal{C}$ there is a canonical map
$$
f^{-1}\SheafHom(X,Y)
\longrightarrow
\SheafHom(f^{-1}X,f^{-1}Y).
$$
It is an isomorphism if, locally on $\Sh(\mathcal{C})$, the object $X$ is
isomorphic to the constant simplicial sheaf associated to a finite
simplicial set.  Here finite means that there are finitely many
nondegenerate simplices; equivalently, the simplicial set has a presentation
as a coequalizer of finite coproducts of standard simplices.
\end{lemma}

\begin{proof}
The degreewise functors are adjoint: under the adjunction
$f^{-1} \dashv f_*$, the compatibility squares for a family of maps
$f^{-1}X_n \to Y_n$ correspond exactly to the compatibility squares for the
adjoint family $X_n \to f_*Y_n$.  Limits and colimits of simplicial objects
are computed degreewise.  The functor $f^{-1}$ preserves all colimits because
it is a left adjoint, by Categories, Lemma
\ref{categories-lemma-adjoint-exact}, and preserves finite limits by the
definition of a morphism of topoi, Sites, Definition
\ref{sites-definition-topos}.  Hence these degreewise functors define the
asserted morphism of topoi.

Apply $f^{-1}$ to the evaluation map and use preservation of finite products:
$$
f^{-1}\SheafHom(X,Y) \times f^{-1}X
\longrightarrow
f^{-1}Y.
$$
Its transpose is the canonical comparison in the statement.

First suppose $X$ is the constant object associated to a finite simplicial
set $E$.  If $E$ is empty, both internal homs are terminal.  Otherwise,
Simplicial Methods, Lemma
\ref{simplicial-lemma-exists-hom-from-simplicial-set-finite}
constructs $\Hom(E,Y)$, and the defining universal property identifies it
with $\SheafHom(E,Y)$.  The proof of that lemma constructs every degree from
finitely many degrees of $Y$ using a finite limit.  Since $f^{-1}$ preserves
finite limits and finite coproducts, it carries this construction to the
identical construction for $f^{-1}Y$.  It also carries the constant object
$E$ to the corresponding constant object.  Thus the comparison is an
isomorphism for $X=E$.

In general choose a covering $\{U_i \to *\}$ such that
$X|_{U_i}$ is constant with finite value $E_i$.  The inverse images
$f^{-1}U_i$ cover the final object.  Pullback to a slice preserves these
internal homs: if $V \to U_i$ is a test object, both adjunctions identify the
relevant maps with maps $V \times X|_{U_i} \to Y|_{U_i}$.  Thus the
restriction of the comparison over $f^{-1}U_i$ is the comparison for the
induced morphism of slice topoi and the constant finite object $E_i$.  It is
an isomorphism by the preceding finite-limit argument.  Isomorphisms of
simplicial sheaves can be checked degreewise and locally.  Therefore the
original comparison is an isomorphism.
\end{proof}
